\documentclass[10pt,a4paper]{amsart}
\usepackage[margin=19mm]{geometry}
\usepackage{amsmath,amssymb,mathtools}
\usepackage[scr=rsfs,cal=euler]{mathalfa}
\usepackage{xfrac}
\usepackage[unicode,hidelinks,bookmarks=false]{hyperref}
\newcommand{\ie}{i.e.\ }
\newcommand{\cf}{cf.\ }
\newcommand{\resp}{respectively\ }
\newcommand{\Subsection}{\subsection}
\providecommand{\nobreakdash}{\nobreak\hbox{-}}
\setlength{\emergencystretch}{2em}
\multlinegap0pt
\usepackage{bm}
\usepackage{bbm}
\usepackage{dsfont}
\usepackage{xspace}
\usepackage{cancel}
\usepackage{mathtools}
\usepackage{amsthm}
\makeatletter
\@ifundefined{nummer}{\newtheorem{nummer}{Nummer}}{}
\@ifundefined{thm}{\newtheorem{thm}[nummer]{Theorem}}{}
\makeatother
\title[The Neighbour Sum Problem on Trees]{The Neighbour Sum Problem on Trees}
\author{Sayan Dutta, Sohom Gupta}
\date{Source: arXiv:2506.23965v2 (CC BY 4.0)}
\begin{document}
\maketitle

\noindent\emph{Source and licence.} The Neighbour Sum Problem on Trees. Version arXiv:2506.23965v2 (CC BY 4.0), distributed under the Creative Commons Attribution 4.0 International licence, as recorded in the arXiv licence metadata for this exact version and shown on the abstract page https://arxiv.org/abs/2506.23965v2. Full rights notice: licenses/arxiv-2506.23965-CC-BY-4.0.txt.\par\smallskip

\noindent\textit{Editorial scope.} Counted unit: the Thm whose source label is \texttt{main} in the article (source0.tex lines 177--179, source label \texttt{main}), delivered together with the article's own complete original proof of it (source0.tex lines 181--223, one \texttt{proof} environment of 43 lines). No mathematics is written, repaired or re-derived: statement and proof are the article's own text line for line. Required context retained uncounted: none. Every definition, display and label of those lines is retained unchanged. Omitted from this package and not redistributed: 1--176, 180--180, 224--575. Nothing inside a retained excerpt is omitted or altered: every retained body is delivered exactly as the article's lines are written, so no omission receipt of any kind accompanies this package. Citations: the retained excerpts carry no citation command at all, so no bibliography entry of the article is transcribed and no reference is redistributed. Reference adaptation: none was needed and none was made; every reference inside the delivered unit resolves inside it, so no unresolved reference or citation is produced. Printed slips kept literal: no printed word, symbol, hypothesis or number of the counted statement or of its proof was altered. Layout and class: the article's own class is \texttt{article} with packages including inputenc, geometry, soul, authblk, tikz, minted, bm, multicol, graphicx, cite, skak, bbm, lmodern, amsmath; the wrapper is the lane's pinned \texttt{amsart} editorial wrapper at 10pt on A4, which restates the theorem-like environments the retained text needs and adds the article's own macro definitions that the retained text needs (none), with extra wrapper packages (bm, bbm, dsfont, xspace, cancel, mathtools). The delivered numbering is the unit's own. Assets: no retained span contains a figure environment, a graphics-inclusion command, a PDF-page inclusion, a table or a TikZ picture; no image file is copied into this package, the package declares no asset dependency and no image file is redistributed with it. Licence: version arXiv:2506.23965v2 of this article is distributed under the Creative Commons Attribution 4.0 International licence, as recorded in the arXiv licence metadata for this exact version and shown on the abstract page https://arxiv.org/abs/2506.23965v2. A full rights notice is delivered at licenses/arxiv-2506.23965-CC-BY-4.0.txt. Characters: every retained excerpt is pure ASCII, so the delivered engine, XeLaTeX with its default Latin Modern text font, reports no missing character.
\begin{thm}\label{main}
    For a finite tree $\mathcal{T}=(\mathcal{V},\:\mathcal{E})$, there exists a pair $(\mathcal{T}, f)$ satisfying the neighbour sum property if and only if there is at least one vertex $v_0\in \mathcal{V}$ such that $S_{1k}(v_0)=1$.
\end{thm}

\begin{proof}
    First, let us assume that there is a solution to the neighbour sum property and call it $f$. We get the corresponding equations for each $v\in \mathcal{V}$. From now on, we will interpret $v$ not just as the vertex, but also as the value assigned to that vertex under the neighbour sum function - in other words, we will denote $f(v)$ as just $v$ wherever appropriate for a neater presentation.
    
    We have $$v_0 = \sum_{r_1} v_1(r_1)$$ for the root vertex. In general, for every vertex in $A_j$, we can write $$v_j(\bar{r}_{j-1}, r_j) = v_{j-1}(\bar{r}_{j-1})+\sum_{r_{j+1}}v_{j+1}(\bar{r}_{j-1}, r_j, r_{j+1})$$ for all $r_j$. 
    
    Suppose the longest path starting at $v_0$ is of length $k$. Equivalently, the nearest neighbours can be indexed up to $k$ with the sets $A_1,\:A_2,...,\:A_k$. For $A_k$, the neighbour sum property becomes $$v_k(\bar{r}_{k-1}, r_k) = v_{k-1}(\bar{r}_{k-1})$$ for all $r_k$.

    Using the equation for a vertex in $A_{k-1}$, we obtain 
    $$v_{k-1}(\bar{r}_{k-2}, r_{k-1}) = v_{k-2}(\bar{r}_{k-2})\left(1-\sum_{r_k} 1\right)^{-1}$$
    and recursively, 
    \begin{align*}
        & v_{k-2}(\bar{r}_{k-3}, r_{k-2}) = v_{k-3}(\bar{r}_{k-3})\left(1-\sum_{r_{k-1}}\left(1-\sum_{r_k} 1\right)^{-1}\right)^{-1} = v_{k-3}(\bar{r}_{k-3})(1-S_{k-1,k}(\bar{r}_{k-2}))^{-1}\\
        &\implies v_{k-3}(\bar{r}_{k-4}, r_{k-3}) = v_{k-4}(\bar{r}_{k-4})\left(1-\sum_{r_{k-2}}\left(1-\sum_{r_{k-1}}\left(1-\sum_{r_k} 1\right)^{-1}\right)^{-1}\right)^{-1} = v_{k-4}(\bar{r}_{k-4})(1-S_{k-2,k}(\bar{r}_{k-3}))^{-1}\\
        &.\\
        &.\\
        &.\\
        &\implies v_1(r_1) = v_0\left(1-\sum_{r_2}\left(\ldots\left(1-\sum_{r_{k-1}}\left(1-\sum_{r_k} 1\right)^{-1}\right)^{-1}\ldots\right)^{-1}\right)^{-1} = v_0(1-S_{2,k}(r_1))^{-1}\\
        &\implies v_0 = \sum_{r_1}v_1(r_1) = v_0\sum_{r_1}\left(1-\sum_{r_2}\left(\ldots\left(1-\sum_{r_{k-1}}\left(1-\sum_{r_k} 1\right)^{-1}\right)^{-1}\ldots\right)^{-1}\right)^{-1} = v_0S_{1,k}
    \end{align*}

    and hence $S_{1,k} = 1$ if $v_0\neq 0$. Since we assumed $f\not\equiv 0$, there exists some $v_0\neq 0$ implying $S_{1,k}(v_0)=1$.

    To show the converse,  we will explicitly construct the function $f$. To do so, let us set some finite $v_0 = a\neq 0$. Assuming the longest path starting from $v_0$ is of length $k$, we can check that the construction $$v_j(\bar{r}_{j-1},r_j)=v_{j-1}(\bar{r}_{j-1})(1-S_{j+1,k}(\bar{r}_j))^{-1}$$ allows $(\mathcal T, f)$ to satisfy the neighbour sum property at all vertices. 
    
    The only thing that remains to be shown that none of these values become infinitely large. This is only possible if for some $j$, $S_{j+1,k}(\bar{r}_j)=1$ and $v_{j-1}(\bar{r}_{j-1})$ is not zero.

    First we show that all $v_1(r_1)$ are finite. We have
    $$v_1(r_1)=v_0(1-S_{2,k}(r_1))^{-1} = a(1-S_{2,k}(r_1))^{-1}$$
    and hence we must have $S_{2,k}(r_1)=1$ for $v_1(r_1)$ to be infinite. However
    $$S_{1,k} = \sum_{r_1}(1-S_{2,k}(r_1))^{-1} = 1$$
    and so neither of the terms in the sum can be infinite. Therefore, $v_1(r_1)$ is finite for every $r_1$.

    
    Let us assume the induction hypothesis that along with $v_0$, all elements in $\bigsqcup_{i=1}^j A_i$ are finite.
    
    Check that
    $$S_{j+2,k}(\bar{r}_{j+1}) = \sum_{r_{j+2}}\left(1-S_{j+3,k}(\bar{r}_{j+2})\right)^{-1}$$
    and since we have formally defined $1/0=\infty$, therefore if any one of the terms in the sum becomes infinite, the entire sum becomes $-\infty$. Suppose $S_{j+3,k}(\bar{r}_j)=1$ for some $\bar{r}_j$. This gives us $$
    v_{j+1}(\bar{r}_{j+1}) = v_{j}(\bar{r}_{j})(1-S_{j+2,k}(\bar{r}_{j+1}))^{-1} = \dfrac{v_{j}(\bar{r}_{j})}{-\infty}=0$$
    since $v_{j}(\bar{r}_j)\in A_j$ and so is finite.

    Therefore, by strong induction, $f$ is finite at all vertices and so if there exists some $v_0\in\mathcal{V}$ such that $S_{1,k}(v_0)=1$, then there exists a pair $(\mathcal{T}, f)$ that satisfies the neighbour sum property. This completes the proof.
\end{proof}

\end{document}
